\documentclass[12pt,a4paper]{article}

% Packages
\usepackage[utf8]{inputenc}
\usepackage[T1]{fontenc}
\usepackage{geometry}
\usepackage{graphicx}
\usepackage{amsmath}
\usepackage{amssymb}
\usepackage{booktabs}
\usepackage{hyperref}
\usepackage{fancyhdr}
\usepackage{titlesec}
\usepackage{enumitem}
\usepackage{xcolor}
\usepackage{listings}
\usepackage{float}
\usepackage{caption}
\usepackage{subcaption}

\geometry{margin=2.5cm, top=3cm, bottom=2.5cm}

% Header with logo
\pagestyle{fancy}
\fancyhf{}
\fancyhead[L]{
\includegraphics[height=1.2cm]{uorlogo.jpg}}
\fancyhead[R]{\small Explainable AI --- MTech Data Science}
\fancyfoot[C]{\thepage}
\renewcommand{\headrulewidth}{0.4pt}
\renewcommand{\footrulewidth}{0.4pt}

% Code listing style
\lstset{
    basicstyle=\ttfamily\small,
    breaklines=true,
    frame=single,
    backgroundcolor=\color{gray!10},
    keywordstyle=\color{blue},
    commentstyle=\color{green!60!black},
    stringstyle=\color{red!70!black},
    language=Python
}

\begin{document}

% Title Page
\begin{titlepage}
    \centering
    
\includegraphics[width=4cm]{uorlogo.jpg}\\[1cm]
    {\Large \textbf{Nirma University}}\\[0.3cm]
    {\large Institute of Technology}\\[0.3cm]
    {\large Department of Computer Science \& Engineering}\\[1.5cm]
    {\LARGE \textbf{Explainable AI}}\\[0.5cm]
    {\Large \textbf{Lab 4: LIME Explanation for Random Forest Classifier}}\\[1.5cm]
    {\large \textbf{(Local Interpretable Model-agnostic Explanations on Iris Dataset)}}\\[2cm]
    \begin{tabular}{ll}
        \textbf{Student Name:} & Soham Dave \\[0.3cm]
        \textbf{Roll Number:} & 25MCD005 \\[0.3cm]
        \textbf{Program:} & MTech Data Science, Semester 2 \\[0.3cm]
        \textbf{Subject:} & Explainable AI \\[0.3cm]
        \textbf{Subject Guide:} & Dr. Nileshkumar Patel \\
    \end{tabular}
    \vfill
    {\large \today}
\end{titlepage}

\tableofcontents
\newpage

% -------------------------------------------------------
\section{Objective}
The objective of this lab is to demonstrate the use of \textbf{LIME (Local Interpretable Model-agnostic Explanations)} for explaining individual predictions of a \textbf{Random Forest classifier}. LIME provides local, per-prediction explanations for any black-box model, making it a powerful tool for post-hoc explainability.

% -------------------------------------------------------
\section{Dataset}

\subsection{Iris Dataset}
\begin{itemize}
    \item \textbf{Source:} \texttt{sklearn.datasets.load\_iris}
    \item \textbf{Task:} Multi-class classification (3 classes)
    \item \textbf{Samples:} 150 (50 per class)
    \item \textbf{Features (4):}
    \begin{enumerate}
        \item Sepal length (cm)
        \item Sepal width (cm)
        \item Petal length (cm)
        \item Petal width (cm)
    \end{enumerate}
    \item \textbf{Target Classes:}
    \begin{enumerate}[start=0]
        \item Setosa
        \item Versicolor
        \item Virginica
    \end{enumerate}
    \item \textbf{Split:} 80/20 train-test (\texttt{random\_state=42})
\end{itemize}

% -------------------------------------------------------
\section{Methodology}

\subsection{Model Training}

\subsubsection{Random Forest Classifier}
A Random Forest is an ensemble of decision trees, each trained on a bootstrap sample with random feature subsets:
\begin{equation}
    \hat{y} = \text{mode}\{h_1(\mathbf{x}), h_2(\mathbf{x}), \ldots, h_B(\mathbf{x})\}
\end{equation}

The model was configured with $B = 100$ estimators and \texttt{random\_state=42}. It achieved \textbf{100\% accuracy} on the test set (30 samples), with perfect precision, recall, and F1-score for all three classes.

\begin{figure}[H]
    \centering
    \includegraphics[width=0.55\textwidth]{images_4/output_1.png}
    \caption{Confusion matrix for the Random Forest classifier on the Iris test set. All 30 test samples are correctly classified (perfect diagonal).}
    \label{fig:confusion_matrix}
\end{figure}

\subsection{LIME Framework}

\subsubsection{Theory}
LIME explains individual predictions of any black-box model $f$ by fitting a simple, interpretable model $g$ locally around the instance of interest $\mathbf{x}$:

\begin{equation}
    \xi(\mathbf{x}) = \arg\min_{g \in G} \; \mathcal{L}(f, g, \pi_{\mathbf{x}}) + \Omega(g)
\end{equation}

where:
\begin{itemize}
    \item $\mathcal{L}(f, g, \pi_{\mathbf{x}})$ measures how well $g$ approximates $f$ in the locality of $\mathbf{x}$, weighted by the proximity kernel $\pi_{\mathbf{x}}$
    \item $\Omega(g)$ is a complexity penalty (e.g., number of non-zero coefficients)
    \item $G$ is the family of interpretable models (typically sparse linear models)
\end{itemize}

\subsubsection{LIME Algorithm Steps}
\begin{enumerate}
    \item \textbf{Perturbation:} Generate $N$ perturbed samples around $\mathbf{x}$ by adding noise to feature values
    \item \textbf{Prediction:} Query the black-box model $f$ on all perturbed samples
    \item \textbf{Weighting:} Assign weights to perturbed samples based on proximity to $\mathbf{x}$ using an exponential kernel: $\pi_{\mathbf{x}}(\mathbf{z}) = \exp\left(-\frac{d(\mathbf{x}, \mathbf{z})^2}{\sigma^2}\right)$
    \item \textbf{Fitting:} Train a weighted linear regression model on the perturbed samples
    \item \textbf{Interpretation:} The linear model's coefficients serve as local feature importance scores
\end{enumerate}

\subsubsection{Implementation}
\begin{enumerate}
    \item Initialized \texttt{LimeTabularExplainer} with training data, feature names, and class names
    \item For each of the 3 classes, selected 5 test samples belonging to that class
    \item Generated LIME explanations for each sample using \texttt{explain\_instance()}
    \item Accumulated feature importance scores across the 5 samples per class
    \item Visualized per-class average feature importance as horizontal bar charts
\end{enumerate}

% -------------------------------------------------------
\section{Results and Analysis}

\begin{figure}[H]
    \centering
    \includegraphics[width=0.75\textwidth]{images_4/output_2.png}
    \caption{LIME feature importance for class Setosa. Petal length and petal width dominate, as Setosa flowers have distinctly small petals.}
    \label{fig:lime_setosa}
\end{figure}

\begin{figure}[H]
    \centering
    \includegraphics[width=0.75\textwidth]{images_4/output_3.png}
    \caption{LIME feature importance for class Versicolor. All four features contribute, with a more balanced distribution reflecting Versicolor's intermediate position between Setosa and Virginica.}
    \label{fig:lime_versicolor}
\end{figure}

\begin{figure}[H]
    \centering
    \includegraphics[width=0.75\textwidth]{images_4/output_4.png}
    \caption{LIME feature importance for class Virginica. Similar to Setosa, petal measurements dominate, but in the opposite direction --- Virginica has the largest petals.}
    \label{fig:lime_virginica}
\end{figure}

\subsection{Interpretation}

\begin{itemize}
    \item \textbf{Setosa:} LIME reveals that petal length and petal width are the dominant features. Setosa flowers have distinctly small petals ($< 2$ cm length), making these features highly discriminative. The model barely uses sepal measurements for Setosa predictions.
    
    \item \textbf{Versicolor:} The feature importance is more distributed across all four features. This is because Versicolor occupies a ``middle ground'' between Setosa and Virginica. LIME shows that no single feature is sufficient --- the model uses a combination of all measurements to identify Versicolor.
    
    \item \textbf{Virginica:} Again, petal features dominate, but with the opposite sign compared to Setosa. Virginica flowers have the largest petals among the three species, making petal length and width the primary distinguishing characteristics.
\end{itemize}

\subsection{Key Observations}
\begin{enumerate}
    \item \textbf{Model-agnostic:} LIME treats the Random Forest as a black box and explains it through local linear approximations. The same technique works for neural networks, SVMs, or any classifier.
    \item \textbf{Local fidelity:} Each explanation is valid only in the neighborhood of the explained instance. Different instances may have different explanations, reflecting the model's adaptive decision logic.
    \item \textbf{Consistency with domain knowledge:} LIME's emphasis on petal features aligns with botanical knowledge --- petal morphology is the primary taxonomic criterion for differentiating Iris species.
\end{enumerate}

% -------------------------------------------------------
\section{Conclusion}
This lab demonstrated that LIME successfully explains individual predictions of a black-box Random Forest:
\begin{itemize}
    \item LIME provides transparent, per-prediction explanations without requiring access to model internals
    \item The explanations are intuitive and consistent with domain knowledge
    \item By aggregating LIME explanations across multiple samples, we can derive class-level interpretations
    \item LIME bridges the gap between model accuracy (100\% in this case) and model interpretability
\end{itemize}

LIME is particularly valuable in high-stakes applications where practitioners need to understand \emph{why} a specific prediction was made, not just \emph{what} was predicted.

\end{document}
